\section{General inputs and principal values}\label{sec:consequences}

We now pass the hard maximal estimate from compactly supported smooth
inputs to all complex $L^3$ inputs.  The main point is to control all
truncations on one set of full measure.  We then prove joint principal-value
convergence and recover the scalar simplex formulation.

Write $p=3/2$ and
\[
 B_{\varepsilon,R}(F,G)(x,y)
 =\int_{\varepsilon<|t|<R}F(x+t,y)G(x,y+t)\,\frac{\dd t}{t}.
\]
Proposition~\ref{prop:compact-maximal} supplies the hard maximal estimate
for compactly supported smooth inputs; denote its constant by $C_M$.

\begin{lemma}[Finite truncations]\label{lem:finite-truncations}
For complex $F,G\in L^3(\R^2)$, there is a set of full measure on which
every finite truncation is absolutely convergent and depends continuously
on $(\varepsilon,R)$ in $0<\varepsilon<R<\infty$.
The resulting family is independent of the representatives of $F,G$ up
to a common null set.  Moreover,
\begin{equation}\label{eq:elementary-bilinear}
 \left\|\int_{\varepsilon<|t|<R}
 |F(x+t,y)G(x,y+t)|\,\frac{\dd t}{|t|}\right\|_p
 \le 2\log(R/\varepsilon)\|F\|_3\|G\|_3.
\end{equation}
\end{lemma}

\begin{proof}
For each fixed $t$, H\"older's inequality and translation invariance give
$\|F(\cdot+t,\cdot)G(\cdot,\cdot+t)\|_p\le\|F\|_3\|G\|_3$.
Minkowski's integral inequality proves
\eqref{eq:elementary-bilinear}.  Apply it to the countable annuli
$1/n<|t|<n$, $n\ge2$.  Outside their combined exceptional null set,
the integrand is absolutely integrable on every compact annulus away
from zero.  Absolute continuity of its integral gives the asserted
endpoint continuity.  If either input changes on a null set, apply the
same bound to the zero $L^3$ difference, and use the same countable
annuli.  This proves independence simultaneously at every endpoint pair.
\end{proof}

Define $B_*(F,G)$ on this full-measure set as the supremum of
$|B_{\varepsilon,R}(F,G)|$ over all finite endpoint pairs, and set it to
zero elsewhere.  Endpoint continuity makes this supremum equal to the
one over positive rational endpoint pairs, so it is measurable.

\begin{proof}[Proof of Theorem~\ref{thm:main}]
Choose compactly supported smooth $f_m\to F$ and $g_m\to G$ in $L^3$.
For a fixed endpoint pair, bilinearity and
\eqref{eq:elementary-bilinear} show that
$B_{\varepsilon,R}(f_m,g_m)\to B_{\varepsilon,R}(F,G)$ in $L^p$.
Fix a finite family of rational endpoint pairs.  The difference of the
two finite maxima of magnitudes is bounded by the sum of the corresponding
truncation differences.  Thus these maxima also converge in $L^p$.
Passing the compact smooth maximal estimate to the limit bounds this
finite maximum by $C_M\|F\|_3\|G\|_3$.  The constant is independent of
the family.  Exhaust all rational endpoint pairs and use monotone
convergence and Lemma~\ref{lem:finite-truncations} to obtain
Theorem~\ref{thm:main}.

We also obtain a useful continuity statement.  For two pairs of inputs,
bilinearity of every truncation gives, almost everywhere,
\begin{equation}\label{eq:maximal-continuity}
 |B_*(F,G)-B_*(f,g)|
 \le B_*(F-f,G)+B_*(f,G-g).
\end{equation}
The bound just established makes the maximal map continuous from
$L^3\times L^3$ to $L^{3/2}$.  Density shows that it is the unique
continuous extension of its compact smooth values.
\end{proof}

\begin{theorem}[Joint bilinear principal values]\label{thm:bilinear-pv}
For all complex $F,G\in L^3(\R^2)$, the joint limit
\[
 B(F,G)=\lim_{\substack{\varepsilon\downarrow0\\R\uparrow\infty}}
            B_{\varepsilon,R}(F,G)
\]
exists almost everywhere and in $L^{3/2}(\R^2)$, with
$\|B(F,G)\|_{3/2}\le C_M\|F\|_3\|G\|_3$.
More precisely,
\begin{equation}\label{eq:maximal-tail-convergence}
 \left\|\sup_{\substack{0<\varepsilon<1/n\\R>n}}
 |B_{\varepsilon,R}(F,G)-B(F,G)|\right\|_{3/2}
 \longrightarrow0.
\end{equation}
The limit is the unique bounded bilinear extension of the
principal-value operator on compactly supported smooth inputs.
\end{theorem}

\begin{proof}
For compactly supported smooth $f,g$ and a fixed $(x,y)$, the function
$h(t)=f(x+t,y)g(x,y+t)$ is smooth and compactly supported.  The truncated
integral is
\[
 \int_\varepsilon^R\frac{h(t)-h(-t)}t\,\dd t.
\]
Its numerator is $O(t)$ at zero and vanishes for large $t$, so it has a
joint limit at every point.

For general inputs, let
\[
 \Omega_n(F,G)=
 \sup_{\substack{0<\varepsilon,\varepsilon'<1/n\\R,R'>n}}
 |B_{\varepsilon,R}(F,G)-B_{\varepsilon',R'}(F,G)|.
\]
These are measurable: by endpoint continuity the supremum may be taken
over rational parameters.  They decrease with $n$ and are bounded by
$2B_*(F,G)$.  Fix smooth approximants $f_m,g_m$ as above.  Outside a
single null set for this countable collection of input pairs,
\[
 \Omega_n(F,G)\le\Omega_n(f_m,g_m)
       +2B_*(F-f_m,G)+2B_*(f_m,G-g_m).
\]
Letting $n\to\infty$, the smooth joint convergence and
Theorem~\ref{thm:main} give
\[
 \left\|\lim_{n\to\infty}\Omega_n(F,G)\right\|_p
 \le 2C_M\bigl(\|F-f_m\|_3\|G\|_3
              +\|f_m\|_3\|G-g_m\|_3\bigr).
\]
The right side tends to zero with $m$.  Hence the limiting oscillation
is zero almost everywhere.  This is the Cauchy criterion for all pairs
of endpoints in the indicated tails, so the joint complex limit exists.
It is measurable by taking a fixed cofinal sequence of endpoint pairs,
and its magnitude is at most $B_*(F,G)$.

Taking $(\varepsilon',R')$ to the limit in the definition of
$\Omega_n$ bounds the supremum in
\eqref{eq:maximal-tail-convergence} by $\Omega_n(F,G)$.
Dominated convergence applies because
$0\le\Omega_n(F,G)\le2B_*(F,G)\in L^p$.
This proves \eqref{eq:maximal-tail-convergence}, and hence joint norm
convergence.  Passing to the norm limit in the bilinear identities for
the truncations proves bilinearity of $B$.  Its bound and uniqueness
follow from Theorem~\ref{thm:main} and density.
\end{proof}

\subsection{Scalar forms and the simplex coordinates}

For complex $F_0,F_1,F_2\in L^3(\R^2)$, define the flat scalar form
\[
 \mathcal L_{\varepsilon,R}(F_0,F_1,F_2)
 =\iiint_{\varepsilon<|t|<R}
 F_0(x,y)F_1(x+t,y)F_2(x,y+t)
 \,\dd x\,\dd y\,\frac{\dd t}{t}.
\]
\begin{corollary}[Flat scalar principal values]\label{cor:scalar-pv}
For complex $F_0,F_1,F_2\in L^3(\R^2)$, every finite flat truncation is
absolutely integrable, and
\[
 \sup_{0<\varepsilon<R<\infty}
 |\mathcal L_{\varepsilon,R}(F_0,F_1,F_2)|
 \le C_M\prod_{v=0}^2\|F_v\|_3.
\]
Its joint scalar principal value exists and equals
$\int F_0 B(F_1,F_2)$.  This limiting form is the unique bounded
trilinear extension of its compact smooth definition.
\end{corollary}

\begin{proof}
H\"older's inequality in $(x,y)$ bounds the absolute integral by
$2\log(R/\varepsilon)\prod_v\|F_v\|_3$.
Thus Fubini identifies the truncated form with
$\int F_0 B_{\varepsilon,R}(F_1,F_2)$.
Theorem~\ref{thm:main} gives its endpoint-uniform bound, and
Theorem~\ref{thm:bilinear-pv} gives joint scalar convergence by pairing
the $L^{3/2}$ limit with $F_0$.  The limiting form is bounded and
complex trilinear.  Density of compact smooth triples proves uniqueness.
\end{proof}

For the simplex form, set
\[
 \Lambda_{\varepsilon,R}(G_0,G_1,G_2)
 =\iiint_{\varepsilon<|a+b+c|<R}
 G_0(a,b)G_1(b,c)G_2(c,a)
 \,\frac{\dd a\,\dd b\,\dd c}{a+b+c}.
\]

The flat and simplex scalar formulations, including equality of their
optimal norm constants, are already related in
\cite[Appendix~B.1]{KovacThieleZorinKranich2015}. We record explicit
norm-preserving maps and track both finite endpoints and their joint limit.

\begin{proposition}[Scalar equivalence]\label{prop:flat-equivalence}
For all complex $G_0,G_1,G_2\in L^3(\R^2)$, every finite simplex
truncation is absolutely integrable, and
\[
 \sup_{0<\varepsilon<R<\infty}
 |\Lambda_{\varepsilon,R}(G_0,G_1,G_2)|
 \le C_M\prod_{v=0}^2\|G_v\|_3.
\]
Its joint scalar principal value exists for every such triple and is
the unique bounded trilinear extension of its compact smooth definition.
The optimal constants in the uniform flat and simplex scalar estimates
are equal.
\end{proposition}

\begin{proof}
Use the determinant-one change of variables
\[
 x=-a,\qquad y=-b,\qquad t=a+b+c,
 \qquad
 a=-x,\quad b=-y,\quad c=t+x+y.
\]
For flat inputs define
\[
 \begin{aligned}
 G_0(u,v)&=F_0(-u,-v),\\
 G_1(u,v)&=F_1(u+v,-u),\\
 G_2(u,v)&=F_2(-v,u+v).
 \end{aligned}
\]
Each pair map has determinant one.  Their inverses are
\[
 F_0(r,s)=G_0(-r,-s),\quad
 F_1(r,s)=G_1(-s,r+s),\quad
 F_2(r,s)=G_2(r+s,-r).
\]
They preserve $L^3$ norms, compact smoothness, and Schwartz class.
Direct substitution identifies the flat and simplex integrands and
their truncation regions.  The absolute integrability, uniform bound,
and joint scalar convergence therefore follow from the flat conclusions.
The inverse maps realize every simplex triple, which proves equality
of the two optimal scalar constants.  Density gives uniqueness of the
bounded trilinear extension.
\end{proof}

Write $C_S$ for the common optimal scalar constant.  Complex $L^p$
duality shows that it is also the optimal constant for the uniform
family of individual bounds
\[
 \|B_{\varepsilon,R}(F,G)\|_{3/2}
 \le C_S\|F\|_3\|G\|_3.
\]
Indeed, for nonzero $U\in L^p$, the function
\[
 H=\frac{\overline U}{|U|}\,
       \frac{|U|^{1/2}}{\|U\|_p^{1/2}},
\]
set to zero where $U=0$, has $L^3$ norm one and satisfies
$\int HU=\|U\|_p$.  Apply this to each bilinear truncation; the zero
case is immediate.  The pointwise maximal optimal constant is at least
$C_S$.
